\begin{figure}[htbp]
\centering
\begin{tikzpicture}[scale=.7,>=stealth]
    % draw vertices and black edges
    \foreach \p in {1, 3, ..., 17}
        \draw (\p,0) -- (\p-1,0);

    % draw the elements of the permutation
    \foreach \p/\l in {1.5/-7, 3.5/3, 5.5/-1, 7.5/4, 9.5/2, 11.5/8, 13.5/-6, 15.5/-5}
        \node at (\p,-1.2) {$\l$};

    % draw the grey edges
    \foreach \b/\e in {0/6, 5/9, 10/3, 4/7, 8/16, 15/14, 13/2, 1/11, 12/17}
        \draw[gray] (\b,0) to  [out=90,in=90] (\e,0);

    % finally, draw the nodes (this must be the last thing to do, if we want them to
    % appear on top of the rest and hence not be overlapped by the edges)
    \foreach \p/\l in {0/0, 1/14, 2/13, 3/5, 4/6, 5/2, 6/1, 7/7, 8/8, 9/3, 10/4, 11/15, 12/16, 13/12, 14/11, 15/10, 16/9, 17/17}
        \node[vertex,draw] at (\p,0) [label=below:$\l$] {};
\end{tikzpicture}
\caption{The breakpoint graph of $\langle -7\ 3\ -1\ 4\ 2\ 8\ -6\ -5\rangle$.} 
\label{fig:breakpoint-graph-example}
\end{figure}
